% !TEX program = pdflatex
% PRINCIPIA PHYSICA pilot --- topic series, paper I of 3.
% Self-contained arXiv-style preprint; no external .bib, no non-standard packages.
\documentclass[12pt,a4paper]{article}

\usepackage[T1]{fontenc}
\usepackage[utf8]{inputenc}
\usepackage[margin=1.15in]{geometry}
\usepackage{amsmath,amssymb,amsthm}
\usepackage{mathrsfs}
\usepackage{enumitem}
\usepackage[hidelinks]{hyperref}

%---------------------------------------------------------------
% Twelve epistemic classes of the pilot contract, plus example,
% counterexample and remark, all on one shared counter so that the
% type of every numbered claim is visible at the point of use.
%---------------------------------------------------------------
\theoremstyle{definition}
\newtheorem{definition}{Definition}[section]
\newtheorem{axiom}[definition]{Axiom}
\newtheorem{assumption}[definition]{Assumption}
\newtheorem{conjecture}[definition]{Conjecture}
\theoremstyle{plain}
\newtheorem{lemma}[definition]{Lemma}
\newtheorem{proposition}[definition]{Proposition}
\newtheorem{theorem}[definition]{Theorem}
\newtheorem{corollary}[definition]{Corollary}
\newtheoremstyle{commentary}{3pt}{3pt}{\small}{}{\bfseries}{.}{.5em}{}
\theoremstyle{commentary}
\newtheorem{interpretation}[definition]{Interpretation}
\newtheorem{empirical}[definition]{Empirical statement}
\newtheorem{established}[definition]{Established physical result}
\newtheorem{newresult}[definition]{New result claimed by this work}
\newtheorem{counterexample}[definition]{Counterexample}
\newtheorem{example}[definition]{Example}
\newtheorem{remark}[definition]{Remark}

% preprint density
\makeatletter
\def\thm@space@setup{\thm@preskip=3pt plus 1pt minus 1pt
                     \thm@postskip=3pt plus 1pt minus 1pt}
\makeatother
\setlength{\abovedisplayskip}{4pt plus 1pt minus 1pt}
\setlength{\belowdisplayskip}{4pt plus 1pt minus 1pt}
\setlength{\abovedisplayshortskip}{2pt}
\setlength{\belowdisplayshortskip}{2pt}
\allowdisplaybreaks

%--------------- shared series notation (identical in all three) --
\newcommand{\R}{\mathbb{R}}
\newcommand{\Cinf}{C^{\infty}}
\newcommand{\FinSymp}{\mathsf{FinSymp}}
\newcommand{\SympRel}{\mathsf{SympRel}}
\newcommand{\CompSys}{\mathsf{CompSys}}
\newcommand{\HamSys}{\mathsf{HamSys}}
\newcommand{\Oc}{\mathcal{O}}
\newcommand{\Mc}{\mathcal{M}}
\newcommand{\Dsf}{\mathsf{D}}
\newcommand{\eps}{\varepsilon}
\newcommand{\pt}{\mathbf{1}}
\newcommand{\graph}{\operatorname{graph}}
\newcommand{\dive}{\operatorname{div}}

%--------------- paper-specific macros (Paper I only) ------------
\newcommand{\Cbb}{\mathbb{C}}
\newcommand{\prm}{\mathrm{pr}}
\newcommand{\Ac}{\mathcal{A}}
\newcommand{\Sp}{\mathit{Sp}}

\title{\bf Compositional Systems in Finite-Dimensional Classical Mechanics:\\
A Typed Reconstruction and Its Failure Modes}
\author{PRINCIPIA PHYSICA Pilot --- Topic 1}
\date{16 August 2026}

\begin{document}
\maketitle
\begin{abstract}
\noindent
We give a typed reconstruction of \emph{system}, \emph{state}, \emph{observable},
\emph{process} and \emph{composition} for finite-dimensional classical mechanics, testing ---
not assuming --- the thesis that \emph{physics is the study of empirically realized
compositional mathematical structures}. The symplectic product is monoidal but not cartesian:
no projection, no diagonal. Centrally, ``these two subsystems interact'' is \emph{not} an
invariant of $(M,\omega,H)$, so compositional structure is additional data fixed empirically
by what an apparatus reads. And the symplectic Hamiltonian systems are not closed under the
operations physicists call composition, the closed class being the Dirac / port-Hamiltonian
one. Every claim carrying epistemic status is typed by one of twelve labels, which share a
counter with the Counterexample and Remark environments.
\end{abstract}

\section{Scope and epistemic conventions}

This paper does not presuppose the pilot thesis, ``\emph{physics is the study of empirically
realized compositional mathematical structures}''; it builds the smallest compositional
apparatus adequate to finite-dimensional classical mechanics, then attacks it. Field theory,
general relativity and quantum theory appear only as \emph{boundary tests}. It is
Paper~I of three, with Paper~II, \emph{Hamiltonian Reconstruction as a Physical Theory
Object} \cite{PrincipiaII}, and Paper~III, \emph{Empirical Realization, Observational
Equivalence, and the Limits of Finite-Dimensional Classical Models} \cite{PrincipiaIII}.
Twelve labels are kept strictly distinct: Definition, Axiom, Assumption, Conjecture, Lemma,
Proposition, Theorem, Corollary, Interpretation, Empirical statement, Established physical
result, New result claimed by this work. Three rules govern them, verbatim from the pilot
contract: (i) an assumption is never converted into a theorem; (ii) every substantial claim is
proved, cited, empirical, or explicitly conjectural; (iii) mathematical equivalence is never
identified with physical equivalence without an empirical argument.\footnote{\emph{Hedges,
retained rather than tidied away.} Residually \emph{uncertain}, flagged not asserted: the
final page of Meyer 1973 (272 vs.\ 273). The journal data for Weinstein's 2009 survey was
subsequently confirmed --- it appeared as Port.\ Math.\ \textbf{67} (2010), 261--278, the form
cited in Papers~III and~IV of this pilot --- but the arXiv preprint is retained here as the
citation of record for the version consulted.}

\section{Systems, states, observables}

\begin{definition}[System, observables, states]\label{def:system}
A \emph{finite-dimensional classical system} is a triple $(M,\omega,H)$: $M$ smooth,
second-countable, $\dim M=2n$; $\omega\in\Omega^{2}(M)$ closed and nondegenerate;
$H\in\Cinf(M)$; write $M^{-}:=(M,-\omega)$. An \emph{observable} is $f\in\Oc:=\Cinf(M,\R)$;
with $\iota_{X_f}\omega=df$ and $\{f,g\}:=\omega(X_f,X_g)$, $(\Oc,\{\cdot,\cdot\})$ is a
Poisson algebra \cite{Arnold1989}. A \emph{point state} is $x\in M$; a \emph{statistical
state} is a Borel probability measure, the point states being its Dirac measures. Unless
stated otherwise $M$ is connected and $X_H$ complete --- neither automatic, since a particle
escaping to infinity in finite time violates completeness.
\end{definition}

\begin{axiom}[Hamiltonian kinematics and dynamics]\label{ax:ham}
The state space of an isolated finite-dimensional classical system is a symplectic manifold,
its observables the smooth real functions on it, and its evolution the flow $\Phi^{t}_{H}$ of
$X_H$, so $\dot f=\{f,H\}$ \cite{AbrahamMarsden1978,Arnold1989}. Labelled an \emph{axiom},
not a theorem: nothing here derives it, and Section~\ref{sec:interconnection} exhibits
ordinary mechanical systems violating it.
\end{axiom}

\begin{assumption}[Observables]\label{asm:obs}
Every $f\in\Cinf(M)$ is a physical observable.
\end{assumption}

\begin{assumption}[Point states]\label{asm:points}
The physically realized states are the point states.
\end{assumption}

Both are false about laboratory practice (Empirical statement~\ref{emp:units}); both remain
assumptions here, and are never converted into theorems.

\section{Processes; partial composition}

\begin{definition}[Processes and the two categories]\label{def:process}
Four strata: (a) the flow $\Phi^{t}_{H}$; (b) a \emph{symplectomorphism} $\psi:M_1\to M_2$,
$\psi^{*}\omega_2=\omega_1$; (c) a \emph{canonical relation} $L:M_1\rightsquigarrow M_2$, a
closed Lagrangian submanifold $L\subset M_1^{-}\times M_2$, so $\dim L=n_1+n_2$ and
$(-\omega_1\boxplus\omega_2)|_L=0$ \cite{Weinstein1982}; (d) a \emph{Dirac relation},
dropping nondegeneracy \cite{Courant1990}. $\FinSymp$ has symplectic manifolds as objects and
symplectomorphisms as morphisms; $\SympRel$ the same objects with canonical relations as
candidate morphisms, composed by $L_2\circ L_1=\pi_{13}(F)$,
$F:=(L_1\times L_2)\cap(M_1\times\Delta_{M_2}\times M_3)$ --- a canonical relation only when
$F$ is a clean intersection and $\pi_{13}|_F$ is proper and injective.
\end{definition}

\begin{established}[$\SympRel$ is not a category]\label{estab:notcat}
Weinstein states the failure himself: canonical relations ``compose well when a transversality
condition is satisfied, but the failure of the most general compositions to be smooth
manifolds means that the canonical relations do not comprise the morphisms of a category''
\cite{Weinstein2009}.
\end{established}

\begin{counterexample}[Composition that is not a submanifold]\label{cex:compose}
Take $M_1=\pt$ (the one-point symplectic manifold), $M_2=\R^{4}$ with
$\omega=dq_1\wedge dp_1+dq_2\wedge dp_2$, $M_3=\R^{2}(q_1,p_1)$. Let
$S(q_1,q_2)=\tfrac{1}{4}q_2^{4}-\tfrac{1}{2}q_1q_2^{2}$ and
$L_1:=\graph(dS)=\{(q_1,q_2,-\tfrac{1}{2}q_2^{2},\,q_2^{3}-q_1q_2)\}$, a closed embedded
Lagrangian submanifold of dimension $2=0+2$. Let $W=\{p_2=0\}$, a hypersurface, hence
coisotropic, with null distribution spanned by $X_{p_2}=\partial/\partial q_2$, and
$R:M_2\rightsquigarrow M_3$ its reduction relation $\{((q_1,q_2,p_1,0),(q_1,p_1))\}$, of
dimension $3=2+1$, with $-\omega_2\boxplus\omega_3$ pulling back to
$-dq_1\wedge dp_1+dq_1\wedge dp_1=0$. Then $L_1\cap W=\{q_2(q_2^{2}-q_1)=0\}$ is the union of
the line $\{q_2=0\}$ and the parabola $\{q_1=q_2^{2}\}$, meeting transversally at the origin
--- \emph{not} a submanifold, so the intersection is not clean. Since
$\partial S/\partial q_1=-\tfrac{1}{2}q_2^{2}$,
\[
 R\circ L_1 = \{(q_1,0):q_1\in\R\}\ \cup\ \{(q_1,-\tfrac{1}{2}q_1): q_1\ge 0\},
\]
a line with a half-ray attached at the origin: not a one-dimensional submanifold of $\R^{2}$,
so $R\circ L_1$ is not a canonical relation. \hfill$\square$
\end{counterexample}

\section{Composition I: monoidal, not cartesian}

The \emph{direct product} of $(M_1,\omega_1,H_1)$ and $(M_2,\omega_2,H_2)$ is
$(M_1\times M_2,\ \prm_1^{*}\omega_1+\prm_2^{*}\omega_2,\ H_1\circ\prm_1+H_2\circ\prm_2)$,
with unit $\pt$.

\begin{proposition}[No projections]\label{prop:noproj}
Let $\dim M=2n$, $\dim N=2m$, $m\ge1$. The graph $\Gamma=\{((x,y),x)\}\subset
(M\times N)^{-}\times M$ of $\prm_M$ is \emph{not} a canonical relation.
\end{proposition}

\begin{proof}
Lagrangian submanifolds of $(M\times N)^{-}\times M$ have dimension
$\bigl((2n+2m)+2n\bigr)/2=2n+m$, while $\dim\Gamma=2n+2m$; these agree only if $m=0$.
Independently, the pullback of $-(\omega_M\boxplus\omega_N)\boxplus\omega_M$ along
$(x,y)\mapsto((x,y),x)$ is $-\omega_M-\omega_N+\omega_M=-\omega_N\neq0$, so $\Gamma$ is not
even isotropic.
\end{proof}

\begin{theorem}[No diagonal]\label{thm:nodiag}
Let $\dim M=2n\ge2$. The graph $\Gamma_\delta=\{(x,(x,x)):x\in M\}$ of the set-theoretic
diagonal is not a canonical relation $M\rightsquigarrow M\times M$.
\end{theorem}

\begin{proof}
$M^{-}\times(M\times M)$ has dimension $6n$, so its Lagrangian submanifolds have dimension
$3n$, whereas $\dim\Gamma_\delta=2n$, and $3n\ne2n$ for $n\ge1$. The stronger,
convention-independent obstruction is isotropy: the pullback of
$(-\omega)\boxplus\omega\boxplus\omega$ along $x\mapsto(x,x,x)$ is
$-\omega+\omega+\omega=\omega$, nondegenerate. So $\Gamma_\delta$ is a \emph{symplectic}
submanifold --- maximally far from isotropic.
\end{proof}

\begin{corollary}[Not cartesian]\label{cor:noncart}
No \emph{copy relation} --- a canonical relation graphing $x\mapsto(x,x)$ --- exists,
$\Gamma_\delta$ being the only candidate. By the characterization of cartesian structure
through natural diagonals and projections (Fox's theorem \cite{Fox1976}; for the monoidal
background see \cite{MacLane1998}),
$(\SympRel,\times,\pt)$ is symmetric monoidal but not cartesian. This is \emph{not} a
classical no-cloning theorem: statistical states can be copied by preparing two systems
identically, and the obstruction is internal to a chosen morphism class. Reading it as a
physical impossibility would be the mathematical-for-physical substitution the contract
forbids.
\end{corollary}

\section{Composition II: decomposition is not intrinsic}\label{sec:decomp}

\begin{lemma}[Non-interaction equals product flow]\label{lem:prodflow}
For a symplectic splitting $M\cong M_1\times M_2$ with $M_1,M_2$ connected,
$H=H_1\circ\prm_1+H_2\circ\prm_2$ up to a constant \emph{iff} $X_H=(X_{H_1},X_{H_2})$.
\end{lemma}
\begin{proof}
Write $X_H=(V(x),W(y))$. Then $dH=\iota_V\omega_1+\iota_W\omega_2$ splits by
bidegree, so $\partial^{2}H/\partial x\,\partial y=0$ and
$H(x,y)=H(x,y_0)+H(x_0,y)-H(x_0,y_0)$; the converse is the same computation backwards.
\end{proof}

\begin{newresult}[Interaction is not an invariant of $(M,\omega,H)$]\label{thm:twodecomp}
Let $M=\R^{4}$ with $\omega=dq_1\wedge dp_1+dq_2\wedge dp_2$ and
\[
 H=\frac{p_1^{2}+p_2^{2}}{2m}+\frac{k}{2}\bigl(q_1^{2}+q_2^{2}\bigr)
 +\frac{c}{2}\bigl(q_1-q_2\bigr)^{2},\qquad m>0,\ k>0,\ c\neq0,\ k+2c>0 .
\]
Then $(M,\omega)$ admits two symplectic direct-product decompositions
\[
 \textbf{A}:\ M\cong \R^{2}(q_1,p_1)\times\R^{2}(q_2,p_2),\qquad
 \textbf{B}:\ M\cong \R^{2}(Q_+,P_+)\times\R^{2}(Q_-,P_-),
\]
$Q_\pm=(q_1\pm q_2)/\sqrt2$, $P_\pm=(p_1\pm p_2)/\sqrt2$, such that $H$ is interacting with
respect to $\textbf{A}$ and non-interacting with respect to $\textbf{B}$. Hence ``is
composite'' and ``its parts interact'' are properties not of $(M,\omega,H)$ but of the pair
$\bigl((M,\omega,H),\text{decomposition}\bigr)$ --- of the system with a choice of which
mutually commutant Poisson subalgebras count as local. \emph{Priority hedge:} the computation
is elementary and we expect it is folklore; our audit located no published statement of the
\emph{classical} case, only the quantum tensor-product results of Zanardi and of Zanardi,
Lidar and Lloyd \cite{Zanardi2001,ZanardiLidarLloyd2004}, outside scope and a boundary test
only. We claim the result while disclaiming novelty of the computation.
\end{newresult}

\begin{proof}
\emph{$\textbf{B}$ is symplectic.} $dQ_+\wedge dP_+ + dQ_-\wedge dP_-
=\tfrac12[(dq_1+dq_2)\wedge(dp_1+dp_2)+(dq_1-dq_2)\wedge(dp_1-dp_2)]=\omega$. In matrix form
the transformation is $T=\mathrm{diag}(S,S)$ on $(q,p)$ with
$S=2^{-1/2}\left(\begin{smallmatrix}1&1\\1&-1\end{smallmatrix}\right)\in O(2)$, and
$T^{\mathsf T}JT=J$ because $S^{\mathsf T}S=I$; so $T\in \Sp(4,\R)$.

\emph{$H$ decouples in $\textbf{B}$.} Orthogonality of $S$ gives
$p_1^{2}+p_2^{2}=P_+^{2}+P_-^{2}$, $q_1^{2}+q_2^{2}=Q_+^{2}+Q_-^{2}$, $q_1-q_2=\sqrt2\,Q_-$,
so $H=\bigl[P_+^{2}/2m+\tfrac{k}{2}Q_+^{2}\bigr]+\bigl[P_-^{2}/2m+\tfrac{k+2c}{2}Q_-^{2}
\bigr]$, uncoupled oscillators of frequencies $\sqrt{k/m}$, $\sqrt{(k+2c)/m}$.

\emph{$H$ does not decouple in $\textbf{A}$.} If $H=a(q_1,p_1)+b(q_2,p_2)$ then
$\partial^{2}H/\partial q_1\partial q_2=0$; but $\partial H/\partial q_1=kq_1+c(q_1-q_2)$, so
$\partial^{2}H/\partial q_1\partial q_2=-c\neq0$.
\end{proof}

The ambiguity is large: $\Sp(4,\R)$ acts transitively on ordered linear symplectic splittings
of $\R^{4}$ with stabilizer $\Sp(2,\R)\times\Sp(2,\R)$, a $4$-parameter family already.

\begin{conjecture}[Generic non-decomposability]\label{conj:generic}
For generic $H$ on a compact symplectic manifold of dimension $\ge4$, no global symplectic
product decomposition makes $H$ non-interacting. We have no proof: by
Lemma~\ref{lem:prodflow} it would force an invariant foliation of codimension $n$, which
KAM-type genericity should obstruct. \emph{Conjecture, not theorem.}
\end{conjecture}

\begin{definition}[Compositional system]\label{def:compsys}
A \emph{compositional system} is $(M,\omega,H,\Ac_1,\Ac_2)$ with $\Ac_1,\Ac_2\subset\Cinf(M)$
Poisson subalgebras that are each other's commutants and generate a dense subalgebra;
equivalently an integrable symplectic-orthogonal splitting $TM=V_1\oplus V_2$. Write
$\CompSys$ for the category, $\HamSys$ for that of systems $(M,\omega,H)$ with $H$-preserving
symplectomorphisms, $U:\CompSys\to\HamSys$ for the forgetful functor.
\end{definition}

\begin{corollary}\label{cor:forget}
$U$ is not injective on isomorphism classes over a fixed $(M,\omega,H)$:
Theorem~\ref{thm:twodecomp} gives two objects with the same image disagreeing on the
interaction predicate. Compositional structure is \emph{strictly additional data}.
\end{corollary}

\begin{empirical}[Units, finite resolution, correlations]\label{emp:units}
$q+p$ is smooth and dimensionally meaningless; every reading has finite precision $\eps>0$;
on a constraint surface only gauge-invariant functions are measured; the decomposition is
fixed in practice by spatial separation. So Assumption~\ref{asm:obs} is false about laboratory
observables, and locality of apparatus is an empirical input, not a consequence of
Axiom~\ref{ax:ham}. Assumption~\ref{asm:points} fails too: a joint measure on
$M_1\times M_2$ is \emph{not} determined by its marginals, and for chaotic systems a point
state is not realizable at precision $\eps$ over a horizon $T$.
\end{empirical}

\begin{newresult}[Status of the founding thesis in this scope]\label{new:thesis}
The strong reading of the pilot thesis --- that compositional structure is determined by the
mathematics of a physical system --- is \emph{false} here, by Theorem~\ref{thm:twodecomp} and
Corollary~\ref{cor:forget}. Only the weak reading survives, ``empirically realized'' carrying
the load: the decomposition datum is what a measurement arrangement supplies (reading $q_1$
needs an apparatus coupled to particle~1 alone, $Q_+$ one coupled to both). That is an
interpretive correspondence, not a theorem, and it can fail --- an apparatus may couple to a
non-integrable family of observables, for which no $\Ac_i$ exists.
\end{newresult}

\section{Composition III: interconnection is not composition}\label{sec:interconnection}

\begin{remark}[A correction to the usual framing]\label{rem:notcomposition}
It is often said that ``composition fails to preserve Hamiltonian structure''. That is
mis-stated: under the cleanness and properness conditions of Definition~\ref{def:process},
composition of canonical relations succeeds. What fails is a family of \emph{different}
operations also called composition --- restriction to a nonintegrable distribution, addition
of a non-Hamiltonian field, restriction to a non-constant-rank constraint set, quotient by a
non-free action --- items (a)--(d).
\end{remark}

\begin{lemma}[Trace and spectrum obstructions]\label{lem:trace}
(i) If $X=X_H$ for some symplectic form near an equilibrium $p$, then in Darboux coordinates
$DX(p)=\left(\begin{smallmatrix}H_{qp}&H_{pp}\\-H_{qq}&-H_{pq}\end{smallmatrix}\right)$, so
$\operatorname{tr}DX(p)=0$, a diffeomorphism invariant. (ii) A linear Hamiltonian field is
$A=JS$, $S=S^{\mathsf T}$, so $\operatorname{spec}(A)=-\operatorname{spec}(A)$. Either
violation certifies non-Hamiltonicity for \emph{every} symplectic form near $p$.
\end{lemma}

\begin{established}[(a) Nonholonomic: almost-Poisson brackets]\label{estab:nonholo}
For a constraint distribution $D\subset TQ$, d'Alembert's principle induces an
\emph{almost-Poisson} bracket on the constrained phase space whose Jacobiator vanishes
identically \emph{iff} $D$ is integrable, i.e.\ iff the constraints are holonomic
\cite{vanderSchaftMaschke1994,KoonMarsden1997}; Koon and Marsden state it as failure of the
reduced two-form to be closed. The vertical rolling disk is the standard witness, its
distribution being bracket-generating.
\end{established}

\begin{proposition}[Chaplygin sleigh: energy conserved, dynamics not Hamiltonian]
\label{prop:sleigh}
The reduced sleigh equations on $\R^{2}(v,\varpi)$ are $\dot v=a\varpi^{2}$,
$\dot\varpi=-\nu v\varpi$, $\nu=ma/(I+ma^{2})>0$, with
$E=\tfrac12 mv^{2}+\tfrac12(I+ma^{2})\varpi^{2}$ conserved:
$\dot E=mav\varpi^{2}-(I+ma^{2})\nu v\varpi^{2}=0$. Yet at each equilibrium $(v_0,0)$,
$v_0\neq0$, the Jacobian
$\left(\begin{smallmatrix}0&0\\0&-\nu v_0\end{smallmatrix}\right)$ has trace $-\nu v_0\ne0$,
so by Lemma~\ref{lem:trace}(i) no symplectic structure near $(v_0,0)$ makes the flow
Hamiltonian. On $\{\pm\varpi>0\}$ the flow \emph{is} Hamiltonian for $E$ with respect to
$\Omega=\bigl((I+ma^{2})/(a\varpi)\bigr)dv\wedge d\varpi$, which blows up across
$\{\varpi=0\}$ and admits no smooth extension.
\end{proposition}

\begin{proposition}[(b) Dissipative coupling leaves the class]\label{prop:damper}
Couple two identical oscillators by a relative-velocity damper: $\dot q_i=p_i/m$,
$\dot p_1=-kq_1-(\gamma/m)(p_1-p_2)$, $\dot p_2=-kq_2-(\gamma/m)(p_2-p_1)$, $\gamma>0$. In the
normal-mode coordinates of Theorem~\ref{thm:twodecomp} the $+$ mode is undamped with spectrum
$\pm i\sqrt{k/m}$, while the $-$ mode obeys $\lambda^{2}+(2\gamma/m)\lambda+k/m=0$, both roots
in the open left half-plane. The spectrum is not closed under $\lambda\mapsto-\lambda$ and the
trace at the origin is $-2\gamma/m\ne0$; by Lemma~\ref{lem:trace} the interconnection is not
Hamiltonian near the origin, although both factors are.
\end{proposition}

\begin{counterexample}[(c) Non-constant rank: no smooth reduced space]\label{cex:rank}
In $\R^{4}$ with $\omega=dq_1\wedge dp_1+dq_2\wedge dp_2$ let $C=\{p_1=0,\ p_2=q_1^{2}/2\}$, a
surface parametrized by $(u,v)\mapsto(u,v,0,u^{2}/2)$. Then
$\iota^{*}\omega=dv\wedge(u\,du)=-u\,du\wedge dv$, of rank $2$ on $\{q_1\ne0\}$ and $0$ on
$\{q_1=0\}$. The characteristic distribution is not a distribution, so $C/\ker$ carries no
smooth manifold structure and the Gotay--Nester--Hinds algorithm \cite{GotayNesterHinds1978}
does not terminate in a reduced symplectic manifold. Equivalently, with
$\Phi=(p_1,\ p_2-q_1^{2}/2)$ one computes $\{\Phi_1,\Phi_2\}=q_1$, so the Dirac matrix has
determinant $q_1^{2}$: the constraints are second class off $\{q_1=0\}$ and first class on it.
The \emph{class} of the constraints is not constant, which is exactly what breaks the Dirac
bracket. \hfill$\square$
\end{counterexample}

\begin{counterexample}[(d) Singular reduction leaves the manifolds]\label{cex:singular}
Reduction of a Hamiltonian $G$-space at a regular value of the momentum map, with the
$G_\mu$-action free and proper, yields a symplectic manifold $J^{-1}(\mu)/G_\mu$
\cite{MarsdenWeinstein1974,Meyer1973}; both hypotheses are essential. Let $M=\Cbb^{2}$ with its
standard form and let $S^{1}$ act with weights $(1,-1)$, the $1{:}{-1}$ resonance. The
momentum map is $J=\tfrac12(|z_1|^{2}-|z_2|^{2})$; the action is free off the origin and has
full isotropy at $0$, so $\mu=0$ is the unique singular value. The invariant ring is generated
by
$\sigma_1=|z_1|^{2}$, $\sigma_2=|z_2|^{2}$, $\sigma_3=\mathrm{Re}(z_1z_2)$,
$\sigma_4=\mathrm{Im}(z_1z_2)$ subject to $\sigma_3^{2}+\sigma_4^{2}=\sigma_1\sigma_2$, and on
$J^{-1}(0)$ we have $\sigma_1=\sigma_2=:s\ge0$, whence
\[
 M_0=J^{-1}(0)/S^{1}\;\cong\;\{(s,\sigma_3,\sigma_4)\in\R^{3}:\
 \sigma_3^{2}+\sigma_4^{2}=s^{2},\ s\ge0\},
\]
a cone, singular at the apex: homeomorphic to $\R^{2}$ but not a smooth symplectic manifold.
It is a symplectic \emph{stratified} space in the sense of Sjamaar and Lerman
\cite{SjamaarLerman1991}. \hfill$\square$
\end{counterexample}

\begin{established}[The compositionally closed class is Dirac, not symplectic]
\label{estab:dirac}
The composition of two Dirac structures is again a Dirac structure
\cite{Courant1990,Cerveraetal2007}, so port-Hamiltonian systems built on them are closed under
power-conserving interconnection \cite{vanderSchaftJeltsema2014}. Symplectic and Poisson
structures lack that closure, as (a)--(d) show; closure requires enlarging the objects (to
Dirac / port-Hamiltonian systems, or for dissipation contact systems
\cite{Bravettietal2017}) or accepting a partial composition. Cervera, van der Schaft and
Ba\~nos note that even Dirac composition can fail in infinite dimensions.
\end{established}

\section{Empirical realization, and a no-go}

\begin{established}[No-interaction theorem: Currie--Jordan--Sudarshan; Leutwyler]
\label{estab:cjs}
Assume: (H1) finitely many point particles, $N\ge2$, on a finite-dimensional phase space with
the canonical Poisson bracket; (H2) a canonical realization of the Poincar\'e algebra by
phase-space functions, $H$ generating time evolution (Dirac's instant form \cite{Dirac1949});
(H3) the particle positions at a common time are half of a canonical system of variables,
transforming under the Euclidean subgroup canonically and geometrically in the same way;
(H4) the world-line condition, that the set of world lines is frame-independent;
(H5) nondegeneracy. Then the only compatible evolution is \emph{free}: all
accelerations vanish. Proved for $N=2$ by Currie, Jordan and Sudarshan \cite{CJS1963}, for
arbitrary $N$ by Leutwyler \cite{Leutwyler1965}.
\end{established}

\begin{interpretation}\label{interp:cjs}
There is no compositional theory of \emph{interacting relativistic subsystems} in this scope.
Each escape damages a different part of the picture: dropping finite-dimensionality imports
field degrees of freedom; dropping canonical positions means subsystem identity is no longer
read off the phase space; dropping the world-line condition makes ``which subsystem is where''
frame-dependent. This interprets a theorem; it does not extend one.
\end{interpretation}

\subsection*{Hidden assumptions, made explicit and separately falsifiable}

\begin{assumption}[Monoidal composition]\label{asm:monoidal}
$A\otimes B$ is determined by $A$ and $B$ up to canonical isomorphism. \emph{Status: refuted
as stated.} For $H=H_1+H_2+V(q_1,q_2)$ the factors do not determine $H$ --- two masses joined
by springs of stiffness $k\ne k'$ are different composites of the same parts --- so no
bifunctor on pairs (phase space, Hamiltonian) produces interacting composites: the interaction
term is irreducibly extra data. Independently, for pair potentials decaying as $r^{-\alpha}$,
$\alpha\le d$, energy is not additive at all, which underlies ensemble inequivalence
\cite{CampaDauxoisRuffo2009} and covers gravitational and unscreened Coulomb systems.
\end{assumption}

\begin{assumption}[Intrinsic subsystems]\label{asm:intrinsic}
A composite determines its decomposition into parts. \emph{Status: refuted} by
Theorem~\ref{thm:twodecomp} and Corollary~\ref{cor:forget}. Adversarially: ask the framework
how many particles $(M,\omega,H)$ has. It cannot say; particle number is not a symplectic
invariant.
\end{assumption}

\begin{assumption}[Isomorphism implies physical equivalence]\label{asm:iso}
Symplectomorphic systems with conjugate Hamiltonians are physically equivalent.
\emph{Status: refuted, and forbidden by the pilot contract.} A pendulum near equilibrium, a
mass--spring and an $LC$ circuit are conjugate at leading order but differ in units,
measurement procedure, domain of validity and failure mode (separatrix, yield point, breakdown
voltage); and since $H$ is defined only up to a constant and a symplectomorphism, ``energy''
is not an invariant of the isomorphism class.
\end{assumption}

\begin{assumption}[Closure under interconnection]\label{asm:closure}
The finite-dimensional symplectic Hamiltonian systems are closed under composition.
\emph{Status: refuted four independent ways} --- Established result~\ref{estab:nonholo} with
Proposition~\ref{prop:sleigh}; Proposition~\ref{prop:damper}; Counterexample~\ref{cex:rank};
Counterexample~\ref{cex:singular} --- the class-level repair being Established
result~\ref{estab:dirac}.
\end{assumption}

\begin{assumption}[Compositional realization]\label{asm:realization}
The interpretation map commutes with composition. \emph{Status: refuted for long-range forces}
(non-additivity, Assumption~\ref{asm:monoidal}); an unexamined isolation/screening assumption
otherwise. This is what the founding thesis most effectively conceals: ``empirically realized
compositional structure'' presumes realization respects composition --- a claim about
interaction ranges, not a mathematical fact.
\end{assumption}

\begin{assumption}[Autonomy]\label{asm:autonomy}
A realized system is an autonomous Hamiltonian system on a symplectic manifold.
\emph{Status: refuted} by Liouville's theorem \cite[\S16]{Arnold1989}: any system with a
volume-contracting attractor --- every real mechanical system with friction --- is not one.
Precision matters: a damped oscillator does admit a time-dependent Lagrangian of
Caldirola--Kanai type, so ``not Hamiltonian'' is false while ``not autonomous symplectic, and
its Hamiltonian is not the energy'' is true.
\end{assumption}

\section{Unresolved tensions, and conclusion}\label{sec:unresolved}

\textbf{U1, totality versus expressiveness:} the direct product is total but expresses only
\emph{non}-interacting composites, while interconnection is expressive but partial. No
finite-dimensional construction known to us is both, and Established result~\ref{estab:notcat}
records that $\SympRel$, the attempt, fails. \textbf{U2, reduction destroys compositeness:}
the two-body Kepler problem, reduced by translations --- mandatory, absolute position not
being observable --- becomes a one-body problem. A composite reduces to a non-composite; we
regard this as the sharpest single tension. \textbf{U3, exact unit, approximate isolation:}
$A\otimes\pt\cong A$ holds on the nose, whereas physical isolation is an approximation that
fails outright for long-range forces (Assumption~\ref{asm:realization}). We know of no
monoidal category coherent only up to $\eps$.

The thesis is neither confirmed nor idle here; it is \emph{corrected}. Physics, in this scope,
is the study of mathematical structures \emph{together with} an empirically supplied
compositional datum that the mathematics does not determine.

\begin{thebibliography}{99}
\small
\bibitem{AbrahamMarsden1978} R.~Abraham, J.~E. Marsden, \emph{Foundations of Mechanics}, 2nd ed., 1978.
\bibitem{Arnold1989} V.~I. Arnold, \emph{Mathematical Methods of Classical Mechanics}, 2nd ed., 1989.
\bibitem{Bravettietal2017} A.~Bravetti, H.~Cruz, D.~Tapias, Contact Hamiltonian mechanics, \emph{Ann. Physics} \textbf{376} (2017), 17--39.
\bibitem{CampaDauxoisRuffo2009} A.~Campa, T.~Dauxois, S.~Ruffo, Statistical mechanics of solvable models with long-range interactions, \emph{Phys. Rep.} \textbf{480} (2009), 57--159.
\bibitem{Cerveraetal2007} J.~Cervera, A.~J. van der Schaft, A.~Ba\~nos, Interconnection of port-Hamiltonian systems and composition of Dirac structures, \emph{Automatica} \textbf{43} (2007), 212--225.
\bibitem{Courant1990} T.~J. Courant, Dirac manifolds, \emph{Trans. Amer. Math. Soc.} \textbf{319} (1990), 631--661.
\bibitem{CJS1963} D.~G. Currie, T.~F. Jordan, E.~C.~G. Sudarshan, Relativistic invariance and Hamiltonian theories of interacting particles, \emph{Rev. Modern Phys.} \textbf{35} (1963), 350--375; erratum ibid.\ \textbf{35}, 1032.
\bibitem{Dirac1949} P.~A.~M. Dirac, Forms of relativistic dynamics, \emph{Rev. Modern Phys.} \textbf{21} (1949), 392--399.
\bibitem{GotayNesterHinds1978} M.~J. Gotay, J.~M. Nester, G.~Hinds, Presymplectic manifolds and the Dirac--Bergmann theory of constraints, \emph{J. Math. Phys.} \textbf{19} (1978), 2388--2399.
\bibitem{KoonMarsden1997} W.~S. Koon, J.~E. Marsden, The Hamiltonian and Lagrangian approaches to the dynamics of nonholonomic systems, \emph{Rep. Math. Phys.} \textbf{40} (1997), 21--62.
\bibitem{Leutwyler1965} H.~Leutwyler, A no-interaction theorem in classical relativistic Hamiltonian particle mechanics, \emph{Nuovo Cimento} \textbf{37} (1965), 556--567.
\bibitem{Fox1976} T.~Fox, Coalgebras and cartesian categories, \emph{Comm. Algebra} \textbf{4} (1976), 665--667.
\bibitem{MacLane1998} S.~Mac~Lane, \emph{Categories for the Working Mathematician}, 2nd ed., 1998.
\bibitem{MarsdenWeinstein1974} J.~E. Marsden, A.~Weinstein, Reduction of symplectic manifolds with symmetry, \emph{Rep. Math. Phys.} \textbf{5} (1974), 121--130.
\bibitem{Meyer1973} K.~R. Meyer, Symmetries and integrals in mechanics, \emph{Dynamical Systems} (1973), 259--272.
\bibitem{SjamaarLerman1991} R.~Sjamaar, E.~Lerman, Stratified symplectic spaces and reduction, \emph{Ann. of Math.} \textbf{134} (1991), 375--422.
\bibitem{vanderSchaftJeltsema2014} A.~J. van der Schaft, D.~Jeltsema, Port-Hamiltonian systems theory: an introductory overview, \emph{Found. Trends Systems Control} \textbf{1} (2014), 173--378.
\bibitem{vanderSchaftMaschke1994} A.~J. van der Schaft, B.~M. Maschke, On the Hamiltonian formulation of nonholonomic mechanical systems, \emph{Rep. Math. Phys.} \textbf{34} (1994), 225--233.
\bibitem{Weinstein1982} A.~Weinstein, The symplectic ``category'', \emph{Diff. Geom. Methods in Math. Phys.} (1982), 45--51.
\bibitem{Weinstein2009} A.~Weinstein, Symplectic categories, arXiv:0911.4133 (2009).
\bibitem{Zanardi2001} P.~Zanardi, Virtual quantum subsystems, \emph{Phys. Rev. Lett.} \textbf{87} (2001), 077901. \emph{Boundary test only.}
\bibitem{ZanardiLidarLloyd2004} P.~Zanardi, D.~A. Lidar, S.~Lloyd, Quantum tensor product structures are observable induced, \emph{Phys. Rev. Lett.} \textbf{92} (2004), 060402. \emph{Boundary test only.}
\bibitem{PrincipiaII} PRINCIPIA PHYSICA Pilot --- Topic 2, \emph{Hamiltonian Reconstruction as a Physical Theory Object}, pilot topic series, Paper~II of 3, 2026.
\bibitem{PrincipiaIII} PRINCIPIA PHYSICA Pilot --- Topic 3, \emph{Empirical Realization, Observational Equivalence, and the Limits of Finite-Dimensional Classical Models}, pilot topic series, Paper~III of 3, 2026.
\end{thebibliography}

\end{document}
